\section{Introduction}

Tool-use agents increasingly depend on external APIs whose schemas do not stay fixed. Fields are renamed, permissions are split, response contracts are restructured, routes are deprecated, and replacements stop being drop-in. These changes are small enough that a system may still look ``compatible'' at a glance, but large enough to flip the correct action. Existing tool-use benchmarks are not built to isolate this regime. They mostly evaluate static interfaces or treat all schema perturbations as if they were the same kind of robustness test.

That framing misses the central distinction. Some changes preserve the intended semantic action: a capable agent should remain invariant. Other changes expose a true capability gap: a capable agent should \emph{not} remain invariant, and should instead choose a different admissible action or decline execution. The problem is therefore not generic lexical robustness, but \emph{canonical semantic actions under real API evolution}.

This paper presents \toolshift, a benchmark and evaluation protocol for that setting. \toolshift{} introduces a deterministic canonical-action evaluator, separates positive transformation families from negative near-orbits, audits an unambiguous core, and governs evaluation with a dev-panel / frozen-blind-panel split. The central contribution is not another generic robustness benchmark, but a precise evaluation target: \emph{schema-visible capability-gap inhibition}.

Our contributions are:
\begin{enumerate}[leftmargin=1.5em]
    \item We formalize canonical-action evaluation for tool use with set-valued admissible outputs, deterministic canonicalization, and explicit positive / negative / impossible splits.
    \item We build an audited real-evolution benchmark with two-panel governance: a 9-family dev panel for diagnosis and a frozen 6-family blind panel for final review, and we connect it to external benchmark assets, public function-calling baselines, and a retrieval-style document baseline.
    \item We show that the strongest retained system is a localized description/contract-aware scaffold baseline, not a learned pair model: it is perfect on the dev panel and remains strongest on frozen blind review, while external public baselines fail on the same blind negative families.
    \item We report a decisive negative result for the original SCC-style ambition: a teacher-distilled canonical-bottleneck run improves dev-panel negative inhibition only slightly and still fails to recover positive retention.
    \item We provide mechanism, protocol, and boundary evidence explaining both the gains and their limits: localized description and contract cues are decisive, negative performance depends materially on set-valued control policy, blind pressure is family-sensitive but not a single-family artifact, and hidden backend shifts remain unsolved.
\end{enumerate}

The final empirical picture is deliberately narrow. The scaffold baseline reaches perfect dev-panel scores and remains strongest on the frozen blind panel with $\CAA=0.917$, $\CAA_\text{clean}=0.958$, $\CAA^{+}=1.000$, $\NOS=0.727$, and $\POC=1.000$. The blind drop is concentrated in negative inhibition, especially on \texttt{slack\_auth}. The strongest learned localizer is more conservative on blind negatives but gives up too much clean and positive retention to replace the scaffold baseline. A final teacher-distilled canonical-bottleneck SCC check also fails to clear our matched-budget go/no-go bar: relative to AugOnly it raises dev-panel \NOS{} only from $0.606$ to $0.636$, while leaving $\CAA^{+}$ and \POC{} at $0.000$.

The remainder of the paper follows this contribution structure. Section~\ref{sec:related} situates \toolshift{} relative to tool-use benchmarks, robustness work, and API-evolution studies. Section~\ref{sec:protocol} defines canonical actions and the positive / negative / impossible split. Section~\ref{sec:benchmark} describes the benchmark and the validation chain. Section~\ref{sec:methods} summarizes the reference methods and the final SCC check. Section~\ref{sec:results} reports dev and blind results. Section~\ref{sec:analysis} gives masking, decision-state, and boundary evidence. Section~\ref{sec:discussion} closes with limitations and the implications for SCC-style learned routes.
